% ═══════════════════════════════════════════════════════════════════════════
%  فصل ۸ — مداخلات اروپایی: از انقلاب شکوهمند تا ناپلئون
% ═══════════════════════════════════════════════════════════════════════════
\chapter{مداخلات اروپایی: از انقلاب شکوهمند تا ناپلئون}
\label{ch:european}

\begin{tcolorbox}[mybox, title={چکیده‌ی فصل}]
این فصل به بررسی مهم‌ترین نمونه‌های نجات‌طلبی و مداخله‌ی خارجی در تاریخ اروپا — از انقلاب شکوهمند انگلستان (۱۶۸۸) تا جنگ‌های ناپلئونی (۱۸۱۵) — می‌پردازد. تمرکز بر دو محور است: (۱) مواردی که نخبگان یا مردم یک کشور اروپایی از قدرت خارجی دعوت کردند؛ (۲) مواردی که قدرت خارجی بدون دعوت مداخله کرد اما ادعای «آزادسازی» داشت. از مدل شش‌وجهی استفاده شده و نتایج با موارد ایرانی مقایسه می‌شود.
\end{tcolorbox}

% ─────────────────────────────────────────────────────────────────────────
\section{انقلاب شکوهمند (۱۶۸۸): موفق‌ترین نمونه‌ی نجات‌طلبی؟}
\label{sec:ch08-glorious}
% ─────────────────────────────────────────────────────────────────────────

\subsection{زمینه و ماهیت}
\label{subsec:ch08-glorious-background}

\concept{انقلاب شکوهمند} (\lr{Glorious Revolution}) رویدادی بود که طی آن پارلمان انگلستان از \histref{ویلیام سوم اورانژ}{استاتهودر هلند} دعوت کرد تا به انگلستان لشکرکشی کند و \enemy{جیمز دوم} (پادشاه کاتولیک) را سرنگون نماید. این رویداد اغلب به‌عنوان \textbf{موفق‌ترین نمونه‌ی تاریخی «مداخله بر اساس دعوت»} معرفی می‌شود.%
\footnote{\lr{Israel, Jonathan. \textit{The Anglo-Dutch Moment: Essays on the Glorious Revolution and Its World Impact}. Cambridge UP, 1991.}}

\subsubsection{چرا دعوت از ویلیام؟}

\begin{enumerate}[nosep, label=\textbf{\arabic*.}]
  \item \textbf{بحران مذهبی:} جیمز دوم کاتولیک بود و تلاش می‌کرد قدرت کاتولیک‌ها را در کشوری عمدتاً پروتستان گسترش دهد.%
  \footnote{\lr{Harris, Tim. \textit{Revolution: The Great Crisis of the British Monarchy, 1685--1720}. Allen Lane, 2006.}}
  \item \textbf{بحران قانون اساسی:} جیمز بدون رضایت پارلمان حکومت می‌کرد و آزادی‌های مدنی را نقض می‌کرد.
  \item \textbf{تولد ولیعهد:} تولد پسر جیمز (ژوئن ۱۶۸۸) یعنی تداوم سلسله‌ی کاتولیک — برای پروتستان‌ها غیرقابل‌تحمل.
  \item \textbf{ویلیام اورانژ:} شوهر مری (دختر پروتستان جیمز) بود و بنابراین ادعای مشروعی به تاج‌وتخت داشت. هم‌زمان استاتهودر هلند بود و منافع ژئوپولیتیکی (مهار فرانسه) داشت.
\end{enumerate}

\subsubsection{«نامه‌ی دعوت» هفت نجیب‌زاده}

در ۳۰ ژوئن ۱۶۸۸، هفت تن از نجیب‌زادگان و سیاستمداران انگلیسی نامه‌ای محرمانه به ویلیام نوشتند و رسماً از او خواستند با نیروی نظامی به انگلستان بیاید:

\begin{histQuote}{نامه‌ی دعوت، ۳۰ ژوئن ۱۶۸۸}
«ما اطمینان داریم که نوزده بیستم مردم این کشور خواهان تغییرند ... اگر جنابعالی به این کشور بیایید، آنقدر قدرتمند خواهید بود که از هرگونه مقاومت جلوگیری شود ... افسران و سربازان بسیاری در ارتش آن‌چنان ناراضی‌اند که به خدمت شما خواهند آمد.»%
\footnote{\lr{Printed in Andrew Browning, ed. \textit{English Historical Documents, 1660--1714}. Eyre and Spottiswoode, 1953, pp.~120--122. The "Immortal Seven" signatories were: Shrewsbury, Devonshire, Danby, Lumley, Compton, Russell, and Sidney.}}
\end{histQuote}

\subsection{تحلیل شش‌وجهی: انقلاب شکوهمند}
\label{subsec:ch08-glorious-analysis}

\begin{tcolorbox}[casebox, title={تحلیل شش‌وجهی: انقلاب شکوهمند (۱۶۸۸)}]
\begin{itemize}[nosep, rightmargin=1em]
  \item \textbf{ساختاری:} بحران مشروعیت جیمز دوم واقعی و عمیق بود. پارلمان — نهاد نمایندگی اصلی — در تضاد بنیادین با شاه قرار داشت. بنابراین «دعوت‌کنندگان» نمایندگی نسبتاً معتبری از جامعه‌ی سیاسی انگلستان داشتند.

  \item \textbf{حقوقی-هنجاری:} مفهوم مدرن حاکمیت ملی هنوز تثبیت نشده بود. از منظر حقوق عرفی انگلستان، پارلمان می‌توانست ادعای مشروعیت کند. ویلیام نیز از طریق ازدواج با مری، ادعای سلسله‌ای داشت.

  \item \textbf{اقتصادی:} هزینه‌های مداخله نسبتاً کم بود (عملاً بدون جنگ جدی). هلند هزینه‌ی لشکرکشی را متحمل شد اما در قبال آن، اتحاد نظامی انگلستان علیه فرانسه را به‌دست آورد — معامله‌ای سودمند.%
  \footnote{\lr{Pincus, Steve. \textit{1688: The First Modern Revolution}. Yale UP, 2009, pp.~213--240.}}

  \item \textbf{نظامی:} ویلیام با حدود ۱۵٬۰۰۰ سرباز و ۴۶۳ کشتی وارد انگلستان شد — یکی از بزرگ‌ترین عملیات‌های آبی‌خاکی تاریخ اروپا. اما \textbf{خونریزی ناچیز بود:} جیمز دوم بدون نبرد جدی فرار کرد و ارتشش متلاشی شد.%
  \footnote{\lr{Jones, J.R. \textit{The Revolution of 1688 in England}. Weidenfeld and Nicolson, 1972, pp.~283--310.}}

  \item \textbf{روان‌شناختی:} دعوت‌کنندگان نسبتاً واقع‌بینانه بودند: ویلیام را نه «ناجی آسمانی» بلکه ابزار عملی برای حل بحران قانون اساسی می‌دیدند. سطح «آرزواندیشی» پایین بود.

  \item \textbf{ژئوپولیتیک:} ویلیام اورانژ منافع ژئوپولیتیک مشخصی داشت: مهار لویی چهاردهم فرانسه. ورود به انگلستان، هلند را از خطر حمله‌ی فرانسه‌ی متحد انگلستان (تحت جیمز) نجات می‌داد.%
  \footnote{\lr{Israel, 1991, pp.~1--43.}}

  \item \textbf{نتیجه (طیف):} \textbf{A}. تصویب \concept{لایحه‌ی حقوق} (\lr{Bill of Rights}, ۱۶۸۹)؛ تثبیت سلطنت مشروطه؛ حفظ استقلال و حاکمیت انگلستان؛ تقویت پارلمان. ویلیام پادشاه شد اما قدرتش محدود به قانون بود.
\end{itemize}
\end{tcolorbox}

\subsection{چرا انقلاب شکوهمند «موفق» بود؟}
\label{subsec:ch08-glorious-success}

\begin{tcolorbox}[successbox, title={عوامل موفقیت انقلاب شکوهمند}]
\begin{enumerate}[nosep, label=\textbf{\arabic*.}]
  \item \textbf{نهاد داخلی قوی:} پارلمان انگلستان قرن‌ها سابقه داشت و توانست پس از مداخله، قدرت را در دست بگیرد — نه مداخله‌گر.

  \item \textbf{استراتژی خروج ذاتی:} ویلیام «پادشاه مشروطه» شد، نه «فرمانده‌ی اشغالگر». ارتش هلندی پس از مدت کوتاهی بازگشت.

  \item \textbf{منافع متقابل واقعی:} ویلیام به اتحاد انگلستان علیه فرانسه نیاز داشت. پارلمان به حذف جیمز. معامله‌ی برد-برد.

  \item \textbf{فقدان مقاومت جدی:} جیمز پایگاه اجتماعی ضعیفی داشت و ارتشش وفادار نبود.

  \item \textbf{مشروعیت سلسله‌ای مداخله‌گر:} ویلیام از طریق همسرش ادعای مشروع به تاج‌وتخت داشت — او «بیگانه‌ی محض» نبود.

  \item \textbf{سطح پایین آرزواندیشی:} دعوت‌کنندگان انتظارات واقع‌بینانه داشتند و توانستند شرایط مداخله را تعریف کنند.
\end{enumerate}
\end{tcolorbox}

\subsection{مقایسه با موارد ایرانی}
\label{subsec:ch08-glorious-comparison}

\begin{table}[htbp]
\centering
\caption{مقایسه‌ی انقلاب شکوهمند با موارد ایرانی}\label{tab:ch08-comparison}
\begin{adjustbox}{max width=\linewidth}
\begin{tabular}{
  >{\raggedleft\arraybackslash\bfseries}p{3cm}
  >{\raggedleft\arraybackslash}p{3.5cm}
  >{\raggedleft\arraybackslash}p{3.5cm}
  >{\raggedleft\arraybackslash}p{3.5cm}}
\toprule
\rowcolor{PrimaryDeep}
\textcolor{white}{عامل} &
\textcolor{white}{انقلاب شکوهمند} &
\textcolor{white}{گرجستان (گئورگیفسک)} &
\textcolor{white}{بحران آذربایجان} \\
\midrule

نهاد داخلی قوی &
\cellcolor{BgSuccess}بله (پارلمان) &
\cellcolor{BgFail}خیر (خاندان ضعیف) &
\cellcolor{BgFail}خیر (فرقه نوپا) \\
\rowcolor{NeutralLight}
مشروعیت مداخله‌گر &
\cellcolor{BgSuccess}بالا (ادعای سلسله‌ای) &
\cellcolor{BgWarm}متوسط (هم‌کیشی) &
\cellcolor{BgFail}پایین (ایدئولوژیک) \\
استراتژی خروج &
\cellcolor{BgSuccess}بله (سلطنت مشروطه) &
\cellcolor{BgFail}خیر (الحاق) &
\cellcolor{BgFail}خیر (عقب‌نشینی ناگهانی) \\
\rowcolor{NeutralLight}
تقارن منافع &
\cellcolor{BgSuccess}بالا (برد-برد) &
\cellcolor{BgFail}پایین (روسیه غالب) &
\cellcolor{BgFail}پایین (شوروی ابزاری) \\
نتیجه &
\cellcolor{BgSuccess}A (موفقیت) &
\cellcolor{BgFail}D–E (الحاق) &
\cellcolor{BgFail}E (فروپاشی) \\

\bottomrule
\end{tabular}
\end{adjustbox}
\end{table}

% ─────────────────────────────────────────────────────────────────────────
\section{جنگ استقلال هلند (۱۵۶۸–۱۶۴۸): دعوت از انگلستان و فرانسه}
\label{sec:ch08-dutch}
% ─────────────────────────────────────────────────────────────────────────

\subsection{زمینه و ماهیت}

جنگ هشتادساله‌ی استقلال هلند از اسپانیا نمونه‌ی مهم دیگری از نجات‌طلبی اروپایی است. شورشیان هلندی بارها از قدرت‌های خارجی — به‌ویژه انگلستان و فرانسه — درخواست کمک نظامی کردند.%
\footnote{\lr{Parker, Geoffrey. \textit{The Dutch Revolt}. Rev. ed. Penguin, 2002.}}

نکته‌ی جالب آنکه هلندی‌ها حتی حاکمیت خود را به بیگانه پیشنهاد دادند: در ۱۵۸۰ دوک آنژو (شاهزاده‌ی فرانسوی) و در ۱۵۸۵ ملکه الیزابت اول انگلستان را به حاکمیت هلند دعوت کردند — هر دو رد کردند (هرچند الیزابت نیروی نظامی محدودی فرستاد).%
\footnote{\lr{Wilson, Charles. \textit{Queen Elizabeth and the Revolt of the Netherlands}. Macmillan, 1970.}}

\subsection{نتیجه و تحلیل}

\begin{tcolorbox}[casebox, title={نتایج نجات‌طلبی هلندی‌ها}]
\begin{itemize}[nosep, rightmargin=1em]
  \item کمک‌های نظامی انگلستان و فرانسه \textbf{محدود و مشروط} بود — نه لشکرکشی تمام‌عیار.
  \item هلندی‌ها عمدتاً با \textbf{نیروی خود} جنگیدند و استقلال را به‌دست آوردند.
  \item استقلال نهایی (صلح وستفالی ۱۶۴۸) نتیجه‌ی ترکیبی از مقاومت داخلی، فرسایش اسپانیا، و تغییر موازنه‌ی اروپایی بود.
  \item \textbf{درس:} نجات‌طلبی زمانی موفق‌تر است که «کمک خارجی» مکمل ظرفیت داخلی باشد نه جایگزین آن.%
  \footnote{\lr{Israel, Jonathan. \textit{The Dutch Republic: Its Rise, Greatness, and Fall, 1477--1806}. Oxford UP, 1995, pp.~155--230.}}
\end{itemize}
\end{tcolorbox}

% ─────────────────────────────────────────────────────────────────────────
\section{لشکرکشی‌های ناپلئون: «آزادسازی» یا امپریالیسم؟}
\label{sec:ch08-napoleon}
% ─────────────────────────────────────────────────────────────────────────

\subsection{ایدئولوژی «آزادسازی ملت‌ها»}

\histref{ناپلئون بناپارت}{۱۷۶۹–۱۸۲۱ م} لشکرکشی‌های خود را با ایدئولوژی \concept{صدور انقلاب فرانسه} و «آزادسازی ملت‌های اروپا از ستم فئودالیسم و استبداد» مشروعیت می‌بخشید. در ابتدا، بخش‌هایی از مردم کشورهای هدف — به‌ویژه روشنفکران، بورژوازی و یهودیان — واقعاً از ارتش فرانسه «استقبال» کردند.%
\footnote{\lr{Broers, Michael. \textit{Europe under Napoleon, 1799--1815}. 2nd ed. Bloomsbury, 2015.}}

\subsection{واکنش‌های محلی: طیفی از استقبال تا مقاومت}

\begin{longtable}{
  >{\raggedleft\arraybackslash\bfseries}p{2.2cm}
  >{\raggedleft\arraybackslash}p{3cm}
  >{\raggedleft\arraybackslash}p{3cm}
  >{\raggedleft\arraybackslash}p{2.5cm}
  >{\raggedleft\arraybackslash}p{2.5cm}}
\caption{واکنش ملت‌های اروپایی به لشکرکشی ناپلئون}\label{tab:ch08-napoleon}\\
\toprule
\rowcolor{PrimaryDeep}
\textcolor{white}{\textbf{کشور/منطقه}} &
\textcolor{white}{\textbf{واکنش اولیه}} &
\textcolor{white}{\textbf{واکنش بلندمدت}} &
\textcolor{white}{\textbf{اصلاحات ناپلئونی}} &
\textcolor{white}{\textbf{میراث}} \\
\midrule
\endfirsthead
\toprule
\rowcolor{PrimaryDeep}
\textcolor{white}{\textbf{کشور/منطقه}} &
\textcolor{white}{\textbf{واکنش اولیه}} &
\textcolor{white}{\textbf{واکنش بلندمدت}} &
\textcolor{white}{\textbf{اصلاحات ناپلئونی}} &
\textcolor{white}{\textbf{میراث}} \\
\midrule
\endhead
\bottomrule
\endlastfoot

ایتالیای شمالی &
استقبال روشنفکران و بورژوازی &
نارضایتی از مالیات و سربازگیری &
قانون مدنی، لغو فئودالیسم &
ملی‌گرایی ایتالیایی%
\footnote{\lr{Grab, Alexander. \textit{Napoleon and the Transformation of Europe}. Palgrave, 2003, pp.~90--120.}} \\
\rowcolor{NeutralLight}
آلمان (راینلند) &
استقبال نسبی &
مقاومت فزاینده &
اصلاحات اداری &
ملی‌گرایی آلمانی \\
اسپانیا &
مقاومت خشونت‌بار فوری (\lr{guerrilla}) &
جنگ استقلال (۱۸۰۸–۱۸۱۴) &
محدود &
«باتلاق اسپانیا»%
\footnote{\lr{Esdaile, Charles. \textit{The Peninsular War: A New History}. Allen Lane, 2002.}} \\
\rowcolor{NeutralLight}
لهستان &
استقبال گسترده &
وفاداری تا پایان &
دوک‌نشین ورشو &
ناکامی و تقسیم مجدد%
\footnote{\lr{Zamoyski, Adam. \textit{1812: Napoleon's Fatal March on Moscow}. Harper, 2004.}} \\
مصر &
بی‌تفاوتی تا مقاومت &
مقاومت مسلحانه &
مدرنیزاسیون محدود &
بیداری ملی‌گرایی مصری \\
\rowcolor{NeutralLight}
روسیه &
مقاومت قاطع &
جنگ میهنی ۱۸۱۲ &
— &
شکست ناپلئون \\

\end{longtable}

\subsection{درس‌های ناپلئونی برای نجات‌طلبی}

\begin{tcolorbox}[defbox, title={درس‌های کلیدی از تجربه‌ی ناپلئونی}]
\begin{enumerate}[nosep, label=\textbf{\arabic*.}]
  \item \textbf{استقبال اولیه ≠ حمایت پایدار:} ملت‌هایی که ابتدا از ناپلئون استقبال کردند، به‌تدریج از مالیات‌های سنگین، سربازگیری اجباری و تحقیر ملی ناراضی شدند.

  \item \textbf{«آزادسازی» اغلب پوشش امپریالیسم است:} ناپلئون ملت‌ها را «آزاد» کرد اما سپس آنها را به دولت‌های دست‌نشانده تبدیل نمود.

  \item \textbf{اصلاحات تحمیلی اما ماندگار:} قانون مدنی ناپلئون (\lr{Code Napoléon})، لغو فئودالیسم و اصلاحات اداری اثرات بلندمدت مثبتی داشت — حتی پس از سقوط ناپلئون. \textbf{پارادوکس: مداخله‌ی امپریالیستی گاه پیامدهای مثبت ناخواسته دارد.}%
  \footnote{\lr{Woolf, Stuart. \textit{Napoleon's Integration of Europe}. Routledge, 1991.}}

  \item \textbf{مقاومت ملی قوی‌ترین عامل شکست مداخله است:} اسپانیا و روسیه نشان دادند که هیچ ارتشی — حتی بزرگ‌ترین ارتش اروپا — نمی‌تواند ملتی مقاوم را تسخیر کند.

  \item \textbf{لهستان: تلخ‌ترین نمونه:} لهستانی‌ها وفادارترین متحدان ناپلئون بودند اما پس از سقوط او، دوباره تقسیم شدند. \concept{وفاداری به مداخله‌گر تضمینی برای پاداش نیست.}
\end{enumerate}
\end{tcolorbox}

% ─────────────────────────────────────────────────────────────────────────
\section{کنگره‌ی وین (۱۸۱۵) و «حق مداخله»}
\label{sec:ch08-vienna}
% ─────────────────────────────────────────────────────────────────────────

\subsection{نظام مترنیخی و سرکوب نجات‌طلبی}

\histref{کنگره‌ی وین}{۱۸۱۴–۱۸۱۵ م} نظم جدیدی در اروپا بنا نهاد که بر اصل \concept{مشروعیت سلطنتی} و \concept{موازنه‌ی قدرت} استوار بود. \thinker{کلمنس فون مترنیخ}، صدراعظم اتریش، سیستمی ایجاد کرد که به موجب آن قدرت‌های بزرگ \textbf{حق داشتند} در هر کشور اروپایی مداخله کنند — نه برای «آزادسازی» بلکه برای \textbf{سرکوب انقلاب}.%
\footnote{\lr{Kissinger, Henry. \textit{A World Restored: Metternich, Castlereagh, and the Problems of Peace, 1812--1822}. Houghton Mifflin, 1957.}}

\begin{legalNote}{پروتکل تروپو (۱۸۲۰)}
«دولت‌هایی که به‌واسطه‌ی انقلاب تغییر یافته‌اند ... از اتحاد اروپایی مستثنی خواهند شد ... متحدان حق دارند با ابزارهای صلح‌آمیز یا در صورت لزوم با زور اسلحه، چنین دولت‌هایی را به مسیر درست بازگردانند.»%
\footnote{\lr{Schroeder, Paul W. \textit{The Transformation of European Politics, 1763--1848}. Oxford UP, 1994, pp.~609--615.}}

این اصل عملاً «نجات‌طلبی معکوس» نظام‌یافته بود: نه مردم از بیگانه علیه حاکم، بلکه حاکمان از بیگانه علیه مردم کمک می‌خواستند.
\end{legalNote}

\subsection{نمونه‌ها: سرکوب با «دعوت»}

\begin{longtable}{
  >{\raggedleft\arraybackslash\bfseries}p{2.5cm}
  >{\raggedleft\arraybackslash}p{2.5cm}
  >{\raggedleft\arraybackslash}p{3cm}
  >{\raggedleft\arraybackslash}p{2.5cm}
  >{\raggedleft\arraybackslash}p{2.5cm}}
\caption{مداخلات «اتحاد مقدس» برای سرکوب انقلاب}\label{tab:ch08-holy-alliance}\\
\toprule
\rowcolor{PrimaryDeep}
\textcolor{white}{\textbf{مورد}} &
\textcolor{white}{\textbf{سال}} &
\textcolor{white}{\textbf{مداخله‌گر}} &
\textcolor{white}{\textbf{دعوت‌کننده}} &
\textcolor{white}{\textbf{نتیجه}} \\
\midrule
\endfirsthead
\toprule
\rowcolor{PrimaryDeep}
\textcolor{white}{\textbf{مورد}} &
\textcolor{white}{\textbf{سال}} &
\textcolor{white}{\textbf{مداخله‌گر}} &
\textcolor{white}{\textbf{دعوت‌کننده}} &
\textcolor{white}{\textbf{نتیجه}} \\
\midrule
\endhead
\bottomrule
\endlastfoot

ناپل &
۱۸۲۱ &
اتریش &
پادشاه ناپل &
سرکوب انقلاب%
\footnote{\lr{Schroeder, 1994, pp.~615--620.}} \\
\rowcolor{NeutralLight}
پیدمونت &
۱۸۲۱ &
اتریش &
پادشاه ساردنی &
سرکوب \\
اسپانیا &
۱۸۲۳ &
فرانسه &
فردیناند هفتم &
بازگشت استبداد مطلق%
\footnote{\lr{Carr, Raymond. \textit{Spain, 1808--1975}. 2nd ed. Oxford UP, 1982, pp.~128--140.}} \\
\rowcolor{NeutralLight}
لهستان &
۱۸۳۱ &
روسیه &
(بدون دعوت — سرکوب شورش) &
خونریزی و الحاق بیشتر \\
مجارستان &
۱۸۴۹ &
روسیه &
اتریش (فرانتس یوزف) &
سرکوب انقلاب مجار%
\footnote{\lr{Rapport, Mike. \textit{1848: Year of Revolution}. Basic Books, 2009, pp.~320--355.}} \\

\end{longtable}

% ─────────────────────────────────────────────────────────────────────────
\section{انقلاب‌های ۱۸۴۸ و مسئله‌ی ناسیونالیسم}
\label{sec:ch08-1848}
% ─────────────────────────────────────────────────────────────────────────

انقلاب‌های ۱۸۴۸ (بهار ملت‌ها) موج بزرگی از نجات‌طلبی ملی-لیبرال بود. ملت‌های تحت سلطه‌ی امپراتوری‌های اتریش، روسیه و عثمانی خواهان آزادی بودند — اما از کجا یاری می‌خواستند؟

\begin{tcolorbox}[defbox, title={پارادوکس ۱۸۴۸: ناسیونالیسم بدون ناجی خارجی}]
بر خلاف انتظار، بیشتر انقلاب‌های ۱۸۴۸ \textbf{بدون کمک خارجی مستقیم} صورت گرفتند. دلایل:
\begin{itemize}[nosep, rightmargin=1em]
  \item \textbf{ایدئولوژی ناسیونالیستی:} ناسیونالیسم سده‌ی نوزدهم تأکید بر «خودبسندگی ملی» داشت — دعوت از بیگانه با ایدئولوژی ناسیونالیستی تناقض داشت.%
  \footnote{\lr{Hobsbawm, Eric. \textit{Nations and Nationalism since 1780}. 2nd ed. Cambridge UP, 1992.}}
  \item \textbf{شکست به‌دلیل تنهایی:} مجارها، لهستانی‌ها و ایتالیایی‌ها هر یک به تنهایی شکست خوردند — زیرا نتوانستند (یا نخواستند) از یکدیگر یا از قدرت بزرگ حمایت بخواهند.
  \item \textbf{استثنا: پیدمونت-ساردنی و فرانسه:} ایتالیایی‌ها در وحدت ایتالیا (۱۸۵۹–۱۸۶۱) از فرانسه (ناپلئون سوم) کمک نظامی خواستند — و به‌قیمت واگذاری نیس و ساووآ به فرانسه.%
  \footnote{\lr{Smith, Denis Mack. \textit{Cavour and Garibaldi 1860}. Cambridge UP, 1985.}}
\end{itemize}
\end{tcolorbox}

% ─────────────────────────────────────────────────────────────────────────
\section{مسئله‌ی شرق: مسیحیان عثمانی و مداخله‌ی اروپا}
\label{sec:ch08-eastern-question}
% ─────────────────────────────────────────────────────────────────────────

\subsection{نجات‌طلبی مسیحیان بالکان}

«مسئله‌ی شرق» (\lr{Eastern Question}) — یعنی سرنوشت اقلیت‌های مسیحی در امپراتوری عثمانی — مهم‌ترین محور نجات‌طلبی در سده‌ی نوزدهم اروپا بود. یونانی‌ها، صرب‌ها، بلغارها و ارمنی‌ها هر یک در مقاطعی از قدرت‌های اروپایی (به‌ویژه روسیه، بریتانیا و فرانسه) خواستار حمایت شدند.%
\footnote{\lr{Anderson, M.S. \textit{The Eastern Question, 1774--1923}. Macmillan, 1966.}}

\subsection{یونان: استقلال با مداخله‌ی خارجی}

\begin{tcolorbox}[casebox, title={نمونه: استقلال یونان (۱۸۲۱–۱۸۲۹)}]
\begin{itemize}[nosep, rightmargin=1em]
  \item \textbf{جنبش داخلی:} شورش یونانیان علیه عثمانی (۱۸۲۱) با رهبری داخلی آغاز شد.
  \item \textbf{نجات‌طلبی:} یونانیان از ابتدا از روسیه (هم‌کیش ارتدوکس)، بریتانیا و فرانسه درخواست کمک کردند.
  \item \textbf{مداخله‌ی نظامی:} نبرد ناوارین (۱۸۲۷) — نیروی دریایی مشترک بریتانیا، فرانسه و روسیه ناوگان عثمانی-مصری را نابود کرد.%
  \footnote{\lr{Brewer, David. \textit{The Greek War of Independence}. Overlook, 2001.}}
  \item \textbf{نتیجه:} استقلال یونان (۱۸۲۹) اما با پادشاه آلمانی‌تبار (اتو باواریایی) — یعنی استقلال ناقص.
  \item \textbf{هزینه:} یونان دهه‌ها تحت نفوذ قدرت‌های بزرگ باقی ماند. «حمایت‌کنندگان» منافع ژئوپولیتیک خود را دنبال کردند.
  \item \textbf{جایگاه طیفی:} \textbf{B} (استقلال اما با وابستگی).
\end{itemize}
\end{tcolorbox}

\subsection{ارمنستان عثمانی: نجات‌طلبی فاجعه‌بار}

\begin{tcolorbox}[enemybox, title={فاجعه‌ی ارمنیان عثمانی: تلخ‌ترین نمونه‌ی نجات‌طلبی}]
نجات‌طلبی ارمنیان عثمانی — یعنی درخواست مکرر از روسیه و قدرت‌های اروپایی برای حمایت — یکی از تلخ‌ترین نمونه‌های تاریخی است:

\begin{enumerate}[nosep, label=\textbf{\arabic*.}]
  \item ماده‌ی ۶۱ معاهده‌ی برلین (۱۸۷۸) «اصلاحات» در استان‌های ارمنی‌نشین عثمانی را تعهد کرد — اما هرگز اجرا نشد.%
  \footnote{\lr{Hovannisian, Richard G. \textit{The Armenian People from Ancient to Modern Times}. Vol.~2. Macmillan, 1997, pp.~203--238.}}
  \item \enemy{سلطان عبدالحمید دوم} نجات‌طلبی ارمنی‌ها را بهانه‌ی سرکوب خونین قرار داد (کشتارهای حمیدی، ۱۸۹۴–۱۸۹۶).
  \item در جنگ جهانی اول، حکومت ترکان جوان با اتهام «خیانت» و «همکاری با روسیه»، ارمنیان را به \enemy{نسل‌کشی} (۱۹۱۵) محکوم کرد — کشتار حدود ۱ تا ۱.۵ میلیون نفر.%
  \footnote{\lr{Akçam, Taner. \textit{A Shameful Act: The Armenian Genocide and the Question of Turkish Responsibility}. Metropolitan, 2006.}}
  \item \textbf{نکته‌ی تلخ:} روسیه — «ناجی» ادعایی — نه قبل و نه حین نسل‌کشی اقدام مؤثری نکرد. ارمنیان شرقی (در ایران و روسیه) نجات یافتند اما ارمنیان غربی (عثمانی) نابود شدند.
\end{enumerate}

\textbf{درس:} نجات‌طلبی نه‌تنها ممکن است بی‌نتیجه بماند، بلکه می‌تواند \textbf{بهانه‌ی سرکوب وحشیانه‌تر} شود. «ادراک حکومت از نجات‌طلبی» (حتی اگر اغراق‌آمیز باشد) کافی است تا فاجعه رقم بخورد — همان‌گونه که در مورد مسیحیان ساسانی نیز دیدیم (فصل ۴).
\end{tcolorbox}

% ─────────────────────────────────────────────────────────────────────────
\section{جدول تطبیقی: نجات‌طلبی اروپایی}
\label{sec:ch08-european-table}
% ─────────────────────────────────────────────────────────────────────────

\begin{landscape}
\begin{table}[htbp]
\centering
\caption{مقایسه‌ی تطبیقی نجات‌طلبی‌های اروپایی}\label{tab:ch08-european-comparative}
\begin{adjustbox}{max width=\linewidth}
\begin{tabular}{
  >{\raggedleft\arraybackslash\bfseries\small}p{2.5cm}
  >{\raggedleft\arraybackslash\small}p{1.5cm}
  >{\raggedleft\arraybackslash\small}p{2cm}
  >{\raggedleft\arraybackslash\small}p{2cm}
  >{\raggedleft\arraybackslash\small}p{2cm}
  >{\raggedleft\arraybackslash\small}p{2.5cm}
  >{\raggedleft\arraybackslash\small}p{2cm}
  >{\raggedleft\arraybackslash\small}p{1.5cm}}
\toprule
\rowcolor{PrimaryDeep}
\textcolor{white}{مورد} &
\textcolor{white}{سال} &
\textcolor{white}{دعوت‌کننده} &
\textcolor{white}{مداخله‌گر} &
\textcolor{white}{نهاد داخلی} &
\textcolor{white}{نتیجه‌ی اصلی} &
\textcolor{white}{پیامد بلندمدت} &
\textcolor{white}{طیف} \\
\midrule

انقلاب شکوهمند &
۱۶۸۸ &
پارلمان &
هلند &
\cellcolor{BgSuccess}قوی &
سلطنت مشروطه &
دموکراسی &
A \\
\rowcolor{NeutralLight}
استقلال هلند &
۱۵۶۸– &
شورشیان &
انگلستان/فرانسه &
\cellcolor{BgWarm}متوسط &
استقلال (۱۶۴۸) &
جمهوری قدرتمند &
A–B \\
ناپلئون-ایتالیا &
۱۷۹۶– &
روشنفکران &
فرانسه &
\cellcolor{BgFail}ضعیف &
دولت دست‌نشانده &
ملی‌گرایی &
C–D \\
\rowcolor{NeutralLight}
ناپلئون-لهستان &
۱۸۰۷ &
وطن‌پرستان &
فرانسه &
\cellcolor{BgFail}ضعیف &
دوک‌نشین موقت &
تقسیم مجدد &
E \\
یونان &
۱۸۲۱ &
شورشیان &
روسیه/بریتانیا/فرانسه &
\cellcolor{BgWarm}در حال شکل‌گیری &
استقلال &
وابستگی &
B \\
\rowcolor{NeutralLight}
ارمنیان عثمانی &
۱۸۷۸– &
نخبگان ارمنی &
روسیه (ادعایی) &
\cellcolor{BgFail}ضعیف &
نسل‌کشی &
فاجعه &
F \\

\bottomrule
\end{tabular}
\end{adjustbox}
\end{table}
\end{landscape}

% ─────────────────────────────────────────────────────────────────────────
\section{جمع‌بندی فصل}
\label{sec:ch08-summary}
% ─────────────────────────────────────────────────────────────────────────

\begin{tcolorbox}[wavebox, title={یافته‌های کلیدی فصل ۸}]
\begin{enumerate}[nosep, label=\textbf{\arabic*.}]
  \item \textbf{تنها یک نمونه‌ی «A» وجود دارد:} انقلاب شکوهمند. عوامل منحصربه‌فرد آن (نهاد پارلمان قوی، مشروعیت سلسله‌ای مداخله‌گر، منافع متقابل، بدون خونریزی) در هیچ مورد دیگری تکرار نشده.

  \item \textbf{کمک خارجی مکمل موفق‌تر از جایگزین است:} هلند (نجات‌طلبی + مقاومت داخلی قوی = موفقیت) در مقابل لهستان (وفاداری به ناپلئون بدون ظرفیت داخلی = شکست).

  \item \textbf{ایدئولوژی «آزادسازی» اغلب پوشش امپریالیسم بوده:} ناپلئون، «اتحاد مقدس»، و پان‌اسلاویسم روسی همگی از ایدئولوژی آزادسازی/حمایت برای پیشبرد منافع ژئوپولیتیک استفاده کردند.

  \item \textbf{نجات‌طلبی می‌تواند بهانه‌ی سرکوب شود:} فاجعه‌ی ارمنیان عثمانی تلخ‌ترین نمونه.

  \item \textbf{نظام مترنیخی نشان داد که «مداخله بر اساس دعوت» ابزار دوطرفه است:} هم حاکمان و هم مردم می‌توانند از بیگانه دعوت کنند — و نتیجه در هر دو حالت، تقویت قدرت مداخله‌گر است.
\end{enumerate}
\end{tcolorbox}

\cleardoublepage